% =====================================================================
%  Appendice B — Codice e dati
%  Inclusa da main.tex dopo \appendix
% =====================================================================

\chapter{Codice e dati}
\label{app:codice}

Il codice di analisi e di generazione, i risultati intermedi, le figure e i
dataset di testo generato sono raccolti in un repository pubblico:

\begin{center}
\footnotesize
\url{https://github.com/SamueleZurlo/Fra-ordine-e-disordine-Proprieta-statistiche-del-linguaggio-generato-da-Large-Language-Model}
\end{center}

\noindent
Il repository contiene i notebook che producono ogni figura e ogni tabella di
questa tesi, i file di risultato in formato aperto e i \num{242} documenti
generati lungo lo sweep di temperatura, con il registro delle condizioni di
produzione di ciascuno: temperatura, seme, parametri di campionamento, numero
di riprese e posizioni dei token di fine sequenza. Il codice è distribuito con
licenza MIT, i dati e le figure con licenza Creative Commons Attribuzione 4.0.

Non vi sono inclusi i corpora di provenienza esterna, che sono materiale
altrui. Al loro posto il file \texttt{FONTI.md} riporta, per ciascun corpus,
la fonte, la data di acquisizione, gli identificatori necessari a riscaricarlo
e un'impronta crittografica dei file impiegati, con la quale chi ripeta
l'acquisizione può stabilire se dispone degli stessi byte. Le voci di
Wikipedia sono riscaricabili a partire dall'elenco dei titoli, anch'esso
incluso; le catture di Grokipedia v0.1 sono permanenti sull'Internet Archive;
il romanzo di riferimento è pubblico dominio.

Fa eccezione la versione attuale di Grokipedia. Come discusso nel
§\ref{sec:paternita}, la piattaforma non pubblica una cronologia delle
revisioni e riscrive le voci senza conservarne lo stato precedente: le catture
impiegate qui sono l'unico testimone di quel testo, e le misure del
Capitolo~\ref{cap:risultati2} che vi si appoggiano possono essere ricontrollate
nei valori derivati ma non rifatte a partire dalla fonte. È un limite della
fonte, non del protocollo, ed è la ragione per cui il §\ref{sec:corpora}
conserva copia locale di tutti i corpora scaricati.
